% MSc dissertation example file, February 2022
%
% Leave one of the documentclass lines uncommented to match your degree.
% You may remove the logo option if it causes problems.
% Do not change any other options.
% \documentclass[logo,msc,adi]{infthesis}     % Adv Design Inf
% \documentclass[logo,msc,ai]{infthesis}      % AI
% \documentclass[logo,msc,cogsci]{infthesis}  % Cognitive Sci
% \documentclass[logo,msc,cs]{infthesis}      % Computer Sci
% \documentclass[logo,msc,cyber]{infthesis}   % Cyber Sec
% \documentclass[logo,msc,datasci]{infthesis} % Data Sci
% \documentclass[logo,msc,di]{infthesis}      % Design Inf
% \documentclass[logo,msc,dsti]{infthesis}    % Data Sci TI
% \documentclass[logo,msc,inf]{infthesis}     % Informatics
\documentclass[logo,msc,ai]{infthesis}           % degree unspecified, do not change except to add your degree
%%%%%%%%%%%%%%%%%%%%%%%%
% Understand any problems and seek approval before assuming it's ok to remove ugcheck.
% \usepackage{msccheck}
\usepackage{msccheck} 

% Include any packages you need below, but don't include any that change the page
% layout or style of the dissertation. By including the ugcheck package above,
% you should catch most accidental changes of page layout though.

\usepackage{microtype} % recommended, but you can remove if it causes problems
\usepackage{xcolor} %remove later this is used for helping edit and draft
\usepackage{caption}
\usepackage{subcaption} %subcaptions for plots and figures next to each other. 
\usepackage{amsmath}  %Check we are allowed
\usepackage{amssymb} %Check we are allowed
\usepackage{amsfonts} %Check we are allowed
\usepackage{booktabs} %Check we are allowed
\usepackage{graphicx} %Check we are allowed
\usepackage{tabularx}
\usepackage{float}
\usepackage{underscore} %Check we are allowed


%%% Colour commands - delete for final draft 
% - Editing Macros (Hierarchy of Importance) -
% 1. FACT CHECK (Critical): Red commands immediate attention for potential inaccuracies.
\newcommand{\fc}[1]{\textcolor{red}{[FactC: #1]}} 
% 2. CITE (High): Orange highlights missing evidence or attribution.
\newcommand{\ct}[1]{\textcolor{orange}{[CITE: #1]}} 
% 3. APPENDIX (Medium): Blue indicates structural moves (shifting text out of the main body).
\newcommand{\app}[1]{\textcolor{blue}{[APP: #1]}} 
% 4. SMALL REWORD (Low): Teal is for minor stylistic edits, mirroring your friend's setup.
\newcommand{\sr}[1]{\textcolor{teal}{[Sreword: #1]}}
% 5. FIGURE (Low): Purple is for figure-related notes (refer to figure, create figure, etc.).
\newcommand{\fig}[1]{\textcolor{purple}{[FIG: #1]}}
% 6. redraft (Low): fix immediadtely marker suggestions and corrections
\newcommand{\fiximm}[1]{\textcolor{brown}{[AI: #1]}}
% 7. LIMITATION (Low): Move to limiations 
\newcommand{\limit}[1]{\textcolor{cyan}{[LIM: #1]}}

%=======================================================================================
\begin{document}
\begin{preliminary}

\title{Physics-informed prediction and spatial attribution of forest top-height growth using repeated LiDAR surveys}
%TODO: DRAFT DIFFERENT VERSIONS OF THE TITLE AND PICK AGAIN AT THE END

\author{Shreeya Gorasia}

\date{\today}
%=======================================================================================

\abstract{TODO: FILL IN}

%TODO: writing: Comeback to this section when the entire report is done 
%=======================================================================================

\maketitle

\newenvironment{ethics}
   {\begin{frontenv}{Research Ethics Approval}{\LARGE}}
   {\end{frontenv}\newpage}

\begin{ethics}
This project was planned in accordance with the Informatics Research
Ethics policy. It did not involve any aspects that required approval
from the Informatics Research Ethics committee.

\standarddeclaration
\end{ethics}

%TODO: writing: Comeback to this section when the entire report is done 
%=======================================================================================

\begin{acknowledgements}
Any acknowledgements go here.
\end{acknowledgements}

%TODO: writing: Comeback to this section when the entire report is done 
%=======================================================================================

\tableofcontents
\end{preliminary}

%TODO: check: Comeback to this when entire report is done and check it and make it fits in one page
%=======================================================================================
\chapter{Introduction}  \label{ch:intro}


\paragraph*{Motivation and problem context}
A Forest Digital Twin is a wider long-term goal for the forestry sector: a dynamic digital representation of forest systems that can support monitoring, forecasting and management planning ~ a shift from static, periodic assessment toward adaptive, data-driven planning. The aim is to combine repeated observations, simulation, and machine-learning(ML) and artificial intelligence (AI) to represent forests as changing systems \cite{baskent2026Review_C1_SS, buonocore2022Proposal_C1_SS}. This is increasingly vital as managers face climate-driven risks—such as extreme storms, destructive winds, and severe summer heatwaves that trigger wildfires and alter growth patterns. At the same time, global forestry is shifting toward continuous cover management, requiring more flexible stand-level models \cite{ForestResearch2009Evidence_C1_SS}. 

Sitka spruce (\textit{Picea sitchensis}) is Britain's primary commercial species, covering approximately 50\% of managed forests \cite{ForestResearch2026Sitka_C1_SS} and contributing roughly \pounds1.1 billion annually to Scotland's economy \cite{ScottishForestry2026Future_C1_SS, mason2011Sitka_C1_SS}. The study focuses on Aberfoyle in Trossachs, Scotland, a Sitka spruce-dominated region with 21 years of repeated LiDAR coverage. Featuring varied topography with exposed, windy ridges to sheltered valleys, providing variation that causes different growth trajectories in different areas of the forest. 


% Height, LiDAR and shared growth relationship
Top height is a key measure of stand development, commonly used in forestry models to estimate downstream metrics like timber volume and site productivity \cite{ForestResearch2026Sitka_C1_SS}, and now can be accurately captured at scale using airborne LiDAR scans \cite{suarezminguez2008Practical_C1_SS}. For most tree species, including Sitka spruce, height development over time is modelled by the Chapman-Richards equation \cite{pommerening2015Methods_C2_CR}, a sigmoid growth curve, with its parameters fixed to represent a stand's maximum height potential and its speed of growth. A fixed parameter curve only captures the average population trend, but local trajectories can deviate from this baseline, due to factors like terrain, climate, soil, wind exposure, and active management all critically constrain or enhance a stand's actual development \cite{socha2023Higher_C1_SS, cameron2015Building_C1_SS}. This creates two related questions: whether local departures from this shared curve are associated with identifiable environmental and management conditions, and whether independent attribution methods agree on which conditions matter.

Foresters have no way to know whether a stand's assigned yield class still holds, or what local conditions are driving it away from that expectation. Whether allowing that relationship to instead vary spatially -- rather than being applied uniformly across an entire region, as current forestry practice including the yield-class system does -- captures real local variation a fixed curve cannot, is a further question this dissertation tests directly. Flexible ML architectures can also capture such relationships more directly than a fixed curve, without assuming a fixed shape at all; deep neural networks (DNNs) and gradient-boosted decision trees (XGBoost), tested here as a supplementary top-height prediction comparison, learn these relationships from observation.

\paragraph*{Research gap and Contribution}
While repeated airborne laser scanning has been applied to model forest height growth and productivity \cite{janiec2023Development_C1_SS}, limited evidence exists on whether local environmental conditions are associated with departures from Sitka spruce's shared height-growth curve, and whether different attribution methods identify consistent conditions. Many current remote-sensing models do not account for spatial dependence between nearby-identical trees, leading to data leakage and inflated performance estimates \cite{Karasiak2022Spatial_C5_EVAL, ploton2020Spatial_C5_EVAL, lynch2025Digital_C4_NN}. Relying on black-box neural networks without explainable AI limits their use for stakeholders who require interpretable insights to support decision-making.

This dissertation focuses on a specific part of the Digital Twin pipeline: using attribution methods to explain local growth deviations. It evaluates whether allowing these relationships to vary across space uncovers patterns that a single global model misses, and tests as a secondary step whether this information improves top-height prediction. Geographically Neural Network-Weighted Regression (GNNWR) serves as the primary model for the spatial attribution analysis. 

\section{Objectives and Research Questions}

Answering all three questions below first required cleaning and combining repeated forest
inventory and LiDAR scan records with multi-source geospatial rasters (terrain, wind, climate
and soil) into one unified modelling dataset (Chapter~\ref{ch:data}).

\begin{table}[!ht]
\small\setlength{\tabcolsep}{4pt}\renewcommand{\arraystretch}{1.0}
\begin{tabularx}{\linewidth}{@{} >{\raggedright\bfseries}p{0.18\linewidth} >{\raggedright\arraybackslash}X @{}}
\noalign{\hrule height 1pt}
\textbf{Research Question} & \textbf{Objective, Sub-objectives \& Own Contribution} \\
\noalign{\hrule height 0.6pt}
\textbf{Q1:} Persistent Curve-Departure Attribution & \textbf{(1.1)} Identify environmental and management conditions tied to a plot's persistent deviation from the shared growth curve. \newline \textbf{(1.2)} Compare attribution across model families (Elastic Net, XGBoost, linear mixed-effects) to evaluate agreement on key conditions. \newline \textbf{Contribution:} Highlights which conditions drive persistent growth departures and maps where model families agree or diverge. \\
\hline
\textbf{Q2:} Spatial Attribution & \textbf{(2.1)} Test whether allowing the environment-deviation relationship to vary spatially (via GNNWR) captures structure missed by global Q1 models. \newline \textbf{(2.2)} Assess plot-level reliability, separating genuine environmental trends from noise and target artefacts. \newline \textbf{Contribution:} Determines if spatial variation enhances attribution and sets boundaries on reading individual extremes as biological signals. \\
\hline
\textbf{Q3:} Prediction \& Physics Guidance & \textbf{(3.1)} Benchmark ML models (LR, RF, XGBoost, DNN) and a physics-informed neural network for top-height prediction using spatial cross-validation. \newline \textbf{(3.2)} Test whether attribution signals from Q1 and Q2 improve predictive accuracy. \newline \textbf{Contribution:} Identifies the top-performing spatial baseline and clarifies the exact performance impact of the physics-informed approach. \\
\noalign{\hrule height 1pt}
\end{tabularx}
\label{tab:aims_objectives}
\end{table}

\paragraph*{Main Finding:}
TODO
%=======================================================================================
%=======================================================================================
\chapter{Related Work and Background}

%=======================================================================================
%=======================================================================================
\chapter{Data}
This chapter describes the dataset used for model training and evaluation, outlining key feature engineering steps and data limitations. Table~\ref{tab:data_summary} summarises all predictor and target variables alongside their sources and temporal coverage. Feature selection rationale, data partitioning, and specific modelling choices are detailed in Chapter~\ref{ch:methodology}.

\section{Study area and cohort coverage}
The dataset covers approximately $16\text{~km} \times 12\text{~km}$ of Aberfoyle and contains 795,645 raw plot observations across six survey years (2002--2023). The records contain airborne laser scanning (ALS)-derived variables, including canopy cover and height percentiles, alongside stand and management fields such as planting year and known felling history (Table~\ref{tab:data-fields}). 
Figure~\ref{fig:data-study-area} shows the study-area topography, the spatial coverage of both cohorts, and an example of standard and boundary-clipped plots.

\begin{figure}[!htbp]
\includegraphics[
    width=1.05\textwidth,
    height=0.48\textheight,
    keepaspectratio
]{figures/fig_data/aberfoyle_area.png}
\caption{\small{Left:Aberfoyle study area. Middle: cohort coverage, Right: plot-level detail}}
\label{fig:data-study-area}
\end{figure}

Due to a LiDAR coverage increase of $37\text{~km}^2$ between 2006 and 2008, two nested cohorts were created (Figure~\ref{fig:cohort_coverage_funnel}): a \emph{4-survey} ($\in \{2008, 2012, 2021, 2023\}$) and a \emph{6-survey} ($\in \{2002, 2006, 2008, 2012, 2021, 2023\}$). The former has broad spatial coverage, while the latter extends temporal trajectories. Sequential filters required complete survey presence, pure Sitka spruce, planting between 1850 and the survey year, ages $0\text{--}200$ years, heights $0\text{--}60\text{~m}$ (ecological max; \cite{mason2011Sitka_C1_SS}), and age $\ge 30$ in 2023 to avoid unreliable young-stand heights (Table~\ref{tab:data-cleaning-funnel} \cite{suarezminguez2008Practical_C1_SS}).

\begin{table}[!htbp]
\centering
\scriptsize
\begin{tabular}{lrrrr}
\toprule
& \multicolumn{2}{c}{4-survey} & \multicolumn{2}{c}{6-survey} \\
\cmidrule(lr){2-3}\cmidrule(lr){4-5}
Stage & Plots & Removed & Plots & Removed \\
\midrule
0. Raw dataset
& 230,359 & -- & 230,359 & -- \\
1. Cohort years present
& 230,345 & 0.01\% & 230,359 & 0.00\% \\
2. Complete survey coverage
& 112,782 & 51.04\% & 18,781 & 91.85\% \\
3. Sitka spruce only
& 72,395 & 35.81\% & 13,897 & 26.01\% \\
4. Valid planting year
& 71,766 & 0.87\% & 13,897 & 0.00\% \\
5. Plausible age (0--200 years)
& 71,766 & 0.00\% & 13,897 & 0.00\% \\
6. Plausible height (0--60 m)
& 71,766 & 0.00\% & 13,897 & 0.00\% \\
7. Mature stands (age $\geq 30$ in 2023)
& \textbf{58,112} & 19.03\%
& \textbf{13,769} & 0.92\% \\
\bottomrule
\end{tabular}
\caption{\small{Plots left after each filter. Percentages show the share removed from the previous row, rather than from the raw dataset. Figure~\ref{fig:data-study-area} shows the final spatial coverage.}}
\label{tab:data-cleaning-funnel}
\end{table}

%------------------------------------------------
\section{Response variable and derived attribution targets}

The base measurement every target in this dissertation is built from is the plot-level 95th
percentile ALS return height above ground (\texttt{elev\_\allowbreak percentile\_\allowbreak
95th}) per survey. Three distinct targets are derived from it: the raw height itself (Q3,
prediction); the mean deviation from the shared Chapman--Richards growth curve, averaged across
a plot's full survey history (Q1, persistent attribution); and the deviation between a plot's
own fitted asymptote and its yield-class-implied asymptote (Q2, spatial attribution). How the
shared curve is fitted and each derived target is constructed is defined in
Chapter~\ref{ch:methodology}. Records also supply planting year, management history, and Forest
Research yield and volume estimates (Table~\ref{tab:data-fields}). To prevent data leakage, any
fields derived directly from the height target were excluded from prediction inputs.


\begin{table}[!htbp]
\scriptsize

\setlength{\tabcolsep}{3pt}
\renewcommand{\arraystretch}{0.95}
\begin{tabular}{
    >{\raggedright\arraybackslash}p{2.6cm}
    >{\raggedright\arraybackslash}p{7.0cm}
    >{\raggedright\arraybackslash}p{3.2cm}
}
\hline
Field & Meaning or derivation & Role in this study \\
\hline
\texttt{elev\_\allowbreak percentile\_\allowbreak 95th} & 95th percentile of ALS return heights above ground, used without correction & Response variable \\
\texttt{Top\_Height95} & $=\texttt{elev\_percentile\_95th}\times1.1$; Forest Research's own corrected height estimate & Retained for reporting but not modelled \\
\texttt{plyr}, \texttt{Age} & Planting year; \texttt{Age} = survey year $-$ \texttt{plyr} & Input feature; used for data cleaning \\
\texttt{yldc} & An existing Forest Research yield-class field, not calculated from this project's LiDAR-derived height response & Excluded from prediction inputs; used to construct Q2's own attribution target (\sr{not just "a comparison variable" -- it is load-bearing for Q2, reworded to say so}) \\
\texttt{whcl} & Ordinal Forest Research windthrow hazard class, 0-6 & Management descriptor \\
\texttt{CanopyCover} & Canopy-cover fraction from the same ALS survey and processing step as the height response & Input feature \\
\texttt{Thin}, \texttt{last\_thinn}, \texttt{next\_thin\_} & Thinning history; \texttt{Thin} is binary (1 if the stand has been thinned, 0 if never thinned); \texttt{last\_thinn}/\texttt{next\_thin\_} record the last and next thinning years & Input feature \\
\texttt{Vol95}, \texttt{area}, \texttt{g1}, \texttt{g3} & Forest Research volume estimate, calculated from height, plot area and canopy cover & Used for disturbance screening \\
\texttt{GYCspec95}, \texttt{p1-p5} & General Yield Class, calculated directly from a plot's own measured height (\texttt{Top\_Height95}) & \texttt{GYCspec95}: excluded from prediction, calculation incorporates the height response, target leakage. \texttt{p1-p5}: excluded from prediction inputs for the same reason, but \texttt{p4}/\texttt{p5} are deliberately used, held fixed, to construct Q2's yield-class-implied asymptote -- excluded and load-bearing are not the same thing \sr{(this dual role was previously collapsed into one blanket "excluded" -- clarified)} \\
\hline
\end{tabular}
\caption{\small{Target, stand and management fields provided by Forest Research.}}
\label{tab:data-fields}
\end{table}

\section{Derived terrain, wind, climate and soil variables}
\label{sec:data-environmental}

Each plot was matched with external terrain, wind, climate and soil variables
(Table~\ref{tab:data-env-sources}). Appendix~\ref{tab:data-env-variables_detail_app} defines
each variable. Every source has a different relationship to time -- some describe something
that does not change (terrain shape), some are aggregated over the interval between two
consecutive surveys (storm activity), some are averaged once across a cohort's whole survey
period, and some are a fixed external climate normal computed over a stated multi-year window
that does not track the survey years at all. This matters because a plot's terrain value never
changes across its survey history, while its climate value might, depending on which of these
four cases applies -- and getting this wrong would let the same terrain/climate signal look like
it has a spurious time trend, or hide a real one. Each row below states which case applies, with
the actual year range where one exists; none of these sources are matched to a plot's individual
survey year (see the note after the table for why not).

\begin{table}[!htbp]
\centering
\scriptsize
\setlength{\tabcolsep}{4pt}
\renewcommand{\arraystretch}{1.05}

\begin{tabular}{
    >{\raggedright\arraybackslash}p{2.0cm}
    >{\raggedright\arraybackslash}p{3.8cm}
    >{\raggedright\arraybackslash}p{3.0cm}
    >{\raggedright\arraybackslash}p{3.7cm}
}
\hline
Data family & Source and resolution & Time basis & Variables used \\
\hline

Terrain
& OS Terrain 50, 50\,m \cite{ordnancesurveyTerrain50_C3_LID}
& Static -- no time dimension
& Elevation, slope, aspect, northness and eastness \\

Soil
& CEH Natural Capital, 50\,m \cite{feeney2026CEHNaturalCapital_C3_LID}; SoilGrids v2.0, 250\,m \cite{poggio2021SoilGrids_C3_LID}
& Static -- no time dimension
& Pedotope, drainage, texture and soil pH \\

Wind exposure
& Global Wind Atlas, 250\,m \cite{davis2023GlobalWindAtlas_C3_LID}
& External average, fixed window: 2006--2015
& Long-term wind speed and exposure \\

Observed wind
& MIDAS Open weather stations \cite{metoffice2019MIDASOpen_C3_LID}
& Per-survey-interval -- aggregated over the gap between two consecutive surveys, not a single
per-year value
& Mean and extreme wind conditions \\

Long-term climate
& CHELSA v2.1, 1\,km \cite{karger2017CHELSA_C3_LID}
& External average, fixed window: 1981--2010 -- a climate normal, not matched to the survey
period at all
& Growing degree days and precipitation \\

Matched climate
& HadUK-Grid, 1\,km \cite{hollis2019HadUKGrid_C3_LID}
& Cohort-averaged -- one value per cohort, averaged across that cohort's own survey years
& Temperature and frost \\

Spatial context
& OS Open Roads and Open Rivers, vector data
\cite{ordnancesurveyOpenRoads_C3_LID,ordnancesurveyOpenRivers_C3_LID}
& Static -- no time dimension
& Distance to roads, rivers and forest edges \\

\hline
\end{tabular}

\caption{\small{External data sources used in the analysis, and each one's actual relationship
to time (not all are the same). Values were matched to plot centres. Full source, resolution,
coverage and extraction details are given in Appendix~\ref{app:data-sources}.}}
\label{tab:data-env-sources}
\end{table}

None of these seven sources are re-matched to a plot's specific survey year -- there is no
``per-survey-year'' row in Table~\ref{tab:data-env-sources} because none of the underlying
external products are available at that resolution for this study area and period; the closest
approximation is MIDAS's per-survey-interval aggregation. The Global Wind Atlas gives long-term
modelled wind, TOPEX measures terrain shelter, and MIDAS records wind at three weather stations.
CHELSA describes long-term climate from 1981 to 2010 \cite{karger2017CHELSA_C3_LID} -- a genuine
climate normal, not a single year's observation, and not updated to track the 2008--2023 survey
period; this is a real temporal mismatch worth keeping in mind wherever CHELSA variables are
interpreted causally. HadUK-Grid temperature and frost values were averaged across each cohort's
own survey years \cite{hollis2019HadUKGrid_C3_LID}, so they differ between the 4-survey and
6-survey cohorts but do not vary within a cohort.


%=======================================================================================
%=======================================================================================
\chapter{Methodology and Experimental Framework} \label{ch:methodology}
\section{Experimental framework}
Q1 attributes a plot's deviation from the baseline growth curve to environmental and management factors. Q2 tests whether spatially varying relationships capture what global models miss. Finally, Q3 evaluates whether these signals improve top-height predictions.

\begin{table}[!ht]
\centering
\footnotesize
\renewcommand{\arraystretch}{1.15}

\begin{tabularx}{\linewidth}{@{} l >{\raggedright\arraybackslash}X >{\raggedright\arraybackslash}X >{\raggedright\arraybackslash}X @{}}
\toprule
Q & Unit and target & Models or comparison & Main evaluation \\
\midrule
Q1 & Plot; mean pooled-CR residual & Elastic Net, XGBoost and two-stage
mixed-effects model & Five-fold spatial block CV; coefficients, SHAP and residual Moran's $I$ \\
Q2 & Plot; difference between fitted and yield-class-implied asymptotes & Elastic Net, XGBoost
and GNNWR & Five-fold spatial block CV; attribution, Moran's $I$ and target audit \\
Q3 & Plot--survey row; 95th-percentile top height, and row-level error reduction against the
fold-matched CR curve & CR, linear regression, random forest, XGBoost, DNN, PINN and PINN-$k$
& Five-fold spatial block CV; random-split and PINN forward-pass diagnostics; CR-error-quartile
correction with permutation control \\
\bottomrule
\end{tabularx}
\caption{\small{Units, targets and evaluation used for each question. Q3's row-level CR-error-quartile
check reuses Q3's own model predictions (conditioned DNN, PINN, PINN-$k$ and XGBoost) rather
than introducing a separate comparison.}}
\label{tab:rq-framework}
\end{table}

\begin{table}[!ht]
\centering
\footnotesize
\renewcommand{\arraystretch}{1.15}
\begin{tabularx}{\linewidth}{@{} >{\raggedright\arraybackslash}p{0.20\linewidth} >{\raggedright\arraybackslash}p{0.24\linewidth} >{\raggedright\arraybackslash}X @{}}
\toprule
Choice & Alternative(s), not used & Why \\
\midrule
CR curve, target/baseline & other growth-curve shapes (Gompertz, logistic) & matches Forestry Commission's own yield-class system; needed for fair Q2 comparison against yield-class benchmark \\
Spatial-block CV & random split; temporal split & compartments held whole, real management units, not random points; random/temporal kept only as separate leakage/robustness checks, not the main evaluation \\
Elastic Net & plain Ridge; plain Lasso & Ridge alone: no sparse, "what matters" answer. Lasso alone: picks one of a correlated pair, drops the other, unstable pick. Elastic Net: sparse and stable together \\
LMM & GAMLSS & real but moderate non-normal residuals (skew $-0.48$, excess kurtosis $0.74$); Python GAMLSS tooling too thin for the size of the problem; would need R or a fragile bridge, not worth the switch \\
GNNWR, not classical GWR & fixed kernel shape and bandwidth & one bandwidth number picked in advance, same everywhere; real density gap by cohort (231+ vs.\ 47 compartments); GNNWR learns the weighting instead of guessing one number for the whole map \\
GNNWR, not geographically weighted random forest & geographically weighted random forest & keeps one interpretable coefficient per variable per place; needed to compare directly against EN's and LMM's coefficients, e.g.\ the topex three-way comparison \\
Fixed rate/shape ($p_4,p_5$), Q2 target & fit all three curve params freely, per plot & avoids overfitting short (4--6 survey) series; cost: forces curve-shape mismatch into the ceiling estimate, drives the extreme values found later \\
Semivariogram-based Moran's $I$ & $k{=}8$ nearest-neighbour & $k{=}8$ only reaches $\sim$150--200m on this dense $\sim$20m grid, just next-door plots; real terrain/wind/climate structure resolves around 1{,}150--1{,}740m; the semivariogram reads the real scale off the data as a curve, instead of guessing a neighbour count \\
\bottomrule
\end{tabularx}
\caption{\small{Main design choices and the alternatives not used.}}
\label{tab:design-choices} %TODO, need to add in to relevant places in the methodology (or some of this can be in backgound)
\end{table}

\paragraph*{Data partitions and leakage control} The primary dataset used the four-survey cohort, while the smaller six-survey cohort served as a sensitivity check with longer observation histories. The unit of analysis varied by research question: Q1 and Q2 used a single aggregated value per plot, whereas Q3 evaluated individual plot-survey observations.

Spatial-block cross-validation was the primary evaluation design. Whole forest compartments were assigned to folds, preventing the same compartment from appearing in both training and testing \cite{roberts2017Cross_C5_EVAL}. Training plots within 60\,m of a validation or test plot were also removed, while all held-out plots were retained. Five folds were balanced approximately by row count. In each rotation, one fold was used for testing, the following fold for validation and the remaining three for training. Each compartment was therefore tested once and predictions from the five test folds were pooled before calculating the overall metrics \cite{ploton2020Spatial_C5_EVAL}.

\begin{figure}[!htbp]
\includegraphics[width=\textwidth]{figures/fig_methodology/method_fig-splits-combined_draft.png}
\caption{Evaluation data-splitting strategies: (A) random plot-level split and (B) compartment-level spatial-block split with a 60\,m buffer.}
\label{fig:splits-combined}
\end{figure}

Continuous variables were scaled using the training data and applied to the validation and test sets. Age was scaled on its own so PINN outputs could be turned into metres per year. Soil types were converted into binary categories. Compartment and block IDs were left out of the models. Two neighbour-height variables were also dropped because they were made using nearby plot heights before the split, which could leak test answers.



\section{Q1: Explaining growth deviations}
Q1 uses the four-survey cohort because the six-survey cohort has too few compartments for spatial analysis. For each plot at survey $t$, the gap from expected height is:
\begin{equation}
r_{it}=H_{it}-H(a_{it}),
\end{equation}
which gives the average error per plot over time:
\begin{equation}
\bar r_i=\frac{1}{T_i}\sum_{t=1}^{T_i}r_{it}.
\label{eq:mean-cr-residual}
\end{equation}
Positive values mark plots that grow faster than the shared baseline curve (defined in Section~\ref{sec:rq1-method}). Stand age is left out of the input features since the baseline curve already uses it.

We test Elastic Net and XGBoost using five-fold spatial cross-validation. Elastic Net tunes its penalty settings on each training fold using an inner loop. XGBoost uses the same tuned settings as Q2 (tree depth 4, learning rate 0.04, up to 500 rounds, with early stopping), chosen after checks showed default settings dropped pooled $R^2$ by 0.07--0.09). Models are scored using pooled test $R^2$.

To account for plots grouped within compartments, we also fit a two-stage mixed model. Stage~1 gets $\bar r_i$ from the baseline curve. Stage~2 fits:
\begin{equation}
\bar r_i=\boldsymbol{\beta}^{\top}\mathbf{x}_i+u_{c[i]}+\varepsilon_i,
\label{eq:mixed-effects}
\end{equation}
where $u_{c[i]}\sim\mathcal{N}(0,\sigma_u^2)$ is a random compartment intercept. Fixed effects are fit using restricted maximum likelihood (REML). To test how much compartment variance the features explain, we switch both models to standard maximum likelihood (ML), because REML numbers cannot be fairly compared when models have different numbers of fixed effects. We report values as mean $\pm$ standard deviation across folds.

Four checks help explain the results: (1) Fit Set~4 without canopy cover. (2) Check correlations between canopy cover and each target. (3)Re-run SHAP values without canopy cover to see if other features take its place. (4) Measure how much slope and TOPEX vary between compartments versus within them.

Comparing Elastic Net, mixed-model coefficients, and XGBoost SHAP values gives three clear views of what drives growth. We also run Moran's $I$ on held-out errors to check for leftover spatial clustering (Section~\ref{sec:eval-robustness}).

% \ai{Proposed additional evaluation, not yet run: every Q1 number above (coefficients, $R^2$, Moran's $I$) reflects a single spatial-block split-seed (42). Q2's GNNWR is already reseeded three times to test sign stability (Section~\ref{sec:q2-method}); repeating that reseed for Q1's Elastic Net and XGBoost, reusing the same five-fold compartment assignments under different seeds, would test whether the slope/TOPEX stability contrast reported in Section~\ref{sec:results-q1-attribution} holds up outside one particular fold assignment, directly answering a gap already flagged in this project's own audit notes rather than leaving it a single-seed number reported as if stable.}
% run Q1 across multiple split seeds to confirm fold stability, matching Q2.}

%------------------------------------------------
\section{Q2: Spatial attribution}

Q1's global models leave significant spatial clustering in their residuals across every feature set (Section~\ref{sec:results-q1-attribution}). This motivates testing if letting the environment-deviation relationship vary across space captures patterns that global models miss. We use Geographically Neural Network-Weighted Regression (GNNWR) \cite{du2020Geographically_C6_GWR} rather than standard GWR because it avoids setting fixed distance kernels or bandwidths, and rather than spatial random forests because it keeps one linear coefficient per feature at each point, matching Q1's models (Section~\ref{sec:background-spatial-eval}). Keeping the model linear is deliberate: it isolates spatial variation as one clean factor, separate from model flexibility, which Q1 already tested via XGBoost.

\limit{GNNWR's coefficients are local, but still linear -- one straight line per variable per place, same as Elastic Net, not XGBoost. Q1 and Q3 both find real non-linear, interaction structure in this data. So an unstable or small local coefficient could mean a real spatial difference, or just a straight line fitting a curved relationship badly in that spot -- the two look the same from GNNWR's output alone. Does not affect the main R2 result (model form is held fixed on purpose, to isolate the spatial question); does limit how much any single local coefficient can be trusted on its own. GNNWR here also rescales one frozen global OLS fit per location (a learned multiplier per variable, not an independent local regression) -- a deliberate stability choice, common where local sample density varies, but a narrower constraint than "locally linear" alone implies.}

Q2 fits a separate height ceiling for each plot while fixing the regional Forestry Commission rate and shape values. For plot $i$ at survey $t$, let:
\begin{align}
g_{it} &= \left(1-e^{-p_{4,i}a_{it}}\right)^{p_{5,i}}, \label{eq:git-curve} \\
\hat y_{\max,i} &= \frac{\sum_t H_{it}g_{it}}{\sum_t g_{it}^2}, \label{eq:plot-ymax-fit} \\
y_{\max,i}^{\mathrm{yldc}} &= \mathrm{yldc}_i p_{2,i}+p_{1,i}+2p_{3,i}, \label{eq:yldc-asymptote} \\
\Delta y_{\max,i} &= \hat y_{\max,i}-y_{\max,i}^{\mathrm{yldc}}. \label{eq:local-ymax-diff}
\end{align}
To construct the attribution target, Q2 fits a plot-specific height ceiling while fixing the regional Forestry Commission rate ($p_{4,i}$) and shape ($p_{5,i}$) parameters to avoid overfitting short series (4--6 surveys). For plot $i$ at survey $t$, the growth term $g_{it}$ is given by Equation~\eqref{eq:git-curve}, and fitting through the origin via ordinary least squares yields the estimated asymptote $\hat y_{\max,i}$ in Equation~\eqref{eq:plot-ymax-fit}. The standard yield-class benchmark $y_{\max,i}^{\mathrm{yldc}}$ is derived from regional lookup parameters in Equation~\eqref{eq:yldc-asymptote}. The final target, $\Delta y_{\max,i}$ (Equation~\eqref{eq:local-ymax-diff}), is the difference between these two ceilings. This formulation isolates growth deviations from environmental features, though fixing $p_4$ and $p_5$ forces any local curve-shape mismatch into $\hat y_{\max,i}$, causing the extreme values analyzed below.

Elastic Net and XGBoost are trained on $\Delta y_{\max}$ using Sets~2--4, following Q1's training setup (Section~\ref{sec:q1-method}). GNNWR is also fit to capture spatial variation [1]:
\begin{equation}
\hat y_i=\beta_0(s_i)+\sum_j\beta_j(s_i)x_{ij}.
\label{eq:gnnwr}
\end{equation}
For location $s_i$, the spatial-weighting network takes scaled distances to all training reference plots and predicts multipliers that adjust a frozen global baseline. All main runs use the full, uncapped training set without fixed bandwidths across the five spatial folds.

\subsection*{Hardware limits and adjustments}
Two GPU limits affected GNNWR. First, taking distance vectors to all $\sim$31,000 training plots initially ran out of GPU memory. An early test capped reference plots proportionally across compartments; resolving a second bottleneck later allowed full, uncapped runs for all headline results. A separate test (Section~\ref{sec:results-q2-spatial}) re-applies this cap to check if sample sparsity explains six-survey instability. Second, GNNWR's built-in tool builds a classical GWR hat matrix for AIC and F-tests, which exceeded memory on large runs. We disabled this matrix for large datasets; this does not affect reported $R^2$ or RMSE, but means AIC values use raw feature counts and classical F-tests are omitted. Neither issue affects Elastic Net or XGBoost.

Feature importance is measured using XGBoost split gain and test-fold SHAP values [2]. A separate check groups features into terrain, wind, soil, climate, and spatial-location categories to measure each group's impact on $R^2$ and residual Moran's $I$ (Set~4, 4-survey cohort).

\subsection*{Target and outlier audit}
To remove bad data, plots are dropped if consecutive surveys match either rule:
\begin{itemize}
\setlength\itemsep{2pt}
\item \textbf{Clearfell:} Top height drops $\ge 25\%$, canopy cover drops $\ge 70\%$, and standing volume drops $\ge 70\%$.
\item \textbf{Measurement error:} Top height drops $\ge 25\%$, while canopy cover drops $\le 10\%$ and volume does not drop.
\end{itemize}
Plots with large height drops that match neither rule are kept.

Extreme fitted ceilings are flagged using standard Tukey $1.5\,\mathrm{IQR}$ fences [3]:
\begin{equation}
\left[Q_1 - 1.5\,\mathrm{IQR},\; Q_3 + 1.5\,\mathrm{IQR}\right],
\label{eq:tukey-fences}
\end{equation}
calculated per cohort, matching past growth-model cleaning steps [4]. Flagged plots remain in main runs, but a sensitivity test checks metrics when excluding them from saved test predictions.

The audit checks if the worst 1\% and 5\% of residual plots recur across Elastic Net, XGBoost, and GNNWR. It tests their residual ranges, GNNWR coefficients, boundary distances, spatial grouping, and wind shelter against normal plots, and re-fits them without 2008 data to see if errors stem from target construction rather than true ecological signals.

%------------------------------------------------
\section{Q3: Prediction and physics guidance}

Q3 tests whether environmental signals improve top-height predictions and evaluates what physics-informed neural networks add. It compares a growth curve baseline, standard tabular models, and neural networks on Set~3 across the five spatial folds, with other sets testing the effect of adding environmental inputs.

\subsection{Chapman--Richards (CR) baseline}
\begin{equation}
H(a) = y_{\max}\left(1 - e^{-ka}\right)^p,
\label{eq:cr}
\end{equation}
where $H(a)$ is height at age $a$, $y_{\max}$ is the height asymptote, $k$ is the growth rate, and $p$ controls curve shape. All three parameters are fit freely by non-linear least squares on each training fold and frozen, keeping test folds fully independent.

\subsection*{Conventional models}
We test three baselines: linear regression, random forest, and XGBoost [1]. For Q3, XGBoost was tuned across 27 parameter settings over five folds using validation $R^2$ (Appendix~\ref{app:hyperparameters}).

\subsection{DNN and physics-informed models}
All neural networks use a three-layer trunk (128 units per layer, Leaky-ReLU) trained with Adam (learning rate $10^{-4}$, weight decay $10^{-5}$, batch size 256, early stopping). The standard DNN minimizes mean squared error ($\mathcal{L}_{\mathrm{data}} = \frac{1}{N}\sum (\hat H_i-H_i)^2$). In the conditioned DNN, environmental variables feed directly into the input. In PINN models, they feed into auxiliary networks that guide the loss functions instead of the prediction trunk.

Equations~\eqref{eq:cr-derivative}--\eqref{eq:total-loss} define the physics-informed loss framework:
\begin{align}
H'(a) &= y_{\max}pk e^{-ka}\left(1-e^{-ka}\right)^{p-1}, \label{eq:cr-derivative} \\
\mathcal{L}_{\mathrm{phys}} &= \frac{1}{N}\sum_{i=1}^{N}\left(\frac{\partial\hat H_i}{\partial a_i}-H'(a_i)\right)^2, \label{eq:physics-loss} \\
\mathcal{L}_{\mathrm{traj}} &= \frac{1}{M}\sum_{j=1}^{M}\left(\frac{\hat H_{j,2}-\hat H_{j,1}}{\Delta a_j}-H'(\bar a_j)\right)^2, \label{eq:trajectory-loss} \\
\mathcal{L} &= \mathcal{L}_{\mathrm{data}} + \lambda_{\mathrm{phys}}\mathcal{L}_{\mathrm{phys}} + \lambda_{\mathrm{traj}}\mathcal{L}_{\mathrm{traj}} + 10^{-5}\sum_{\theta}|\theta|. \label{eq:total-loss}
\end{align}
Equation~\eqref{eq:cr-derivative} gives the theoretical growth rate. Equation~\eqref{eq:physics-loss} penalizes gaps between this rate and network derivatives ($\partial\hat H_i/\partial a_i$). Equation~\eqref{eq:trajectory-loss} penalizes differences across consecutive surveys from the same plot at midpoint age $\bar a_j$. The total loss in Equation~\eqref{eq:total-loss} balances these terms ($\lambda_{\mathrm{phys}}=\lambda_{\mathrm{traj}}=1$ by default; 0 disables physics). 

For conditioned PINNs, auxiliary networks adjust plot parameters:
\begin{equation}
y_{\max}^{(i)} = y_{\max}+f_y(\mathbf{x}_i), \quad k^{(i)} = k\exp\!\left(f_k(\mathbf{x}_i)\right).
\label{eq:param-adjust}
\end{equation}
\limit{In the evaluated implementation, $y_{\max}^{(i)}$ and $k^{(i)}$ guide the loss terms rather than passing directly into the prediction trunk; see Section~\ref{sec:rq1-pinn-limitation}.}

\subsection*{Row-level curve correction}
\label{sec:rq2a-method}
To test whether environmental inputs improve single plot--survey predictions, we measure absolute error reduction against the CR baseline:
\begin{equation}
d_i = \left|H_i-\hat H_i^{\mathrm{CR}}\right| - \left|H_i-\hat H_i^{\mathrm{model}}\right|.
\label{eq:residual-reduction}
\end{equation}
Positive $d_i$ indicates improvement. Test rows are split into baseline-error quartiles (Q1 best fit, Q4 worst fit) to see where gains concentrate. Two Set~3 XGBoost controls isolate this effect: removing environmental variables, and refitting from scratch on permuted environmental data. Spatial clustering of $d_i$ is checked with Moran's $I$.


\section{}
EVALUATION : BASICS R2 RMSE ETC 

moran i 
  - Weights: binary distance-band (within threshold = neighbour, else not), row-standardised — not inverse-distance, not a kernel
  - Global Moran's I only, no local variant (LISA) — opposing local clusters could cancel out and go undetected
  - Detects spatial pattern, not its cause — separate diagnostics (outlier audit, category attribution) needed to explain why
  - Lower Moran's I is not the same claim as lower residual variance — keep the two metrics distinct in any comparison
  - Large datasets subsampled to 5,000 rows (fixed seed) before computing, for feasibility

%------------------------------------------------
\section{Evaluation and robustness}
\label{sec:eval-robustness}

For residual $e_i=H_i-\hat H_i$, the main metrics are
\begin{align}
\mathrm{RMSE}&=\sqrt{\frac{1}{N}\sum_i e_i^2}, &
\mathrm{MAE}&=\frac{1}{N}\sum_i|e_i|, \\
R^2&=1-\frac{\sum_i e_i^2}{\sum_i(H_i-\bar H)^2}, &
\mathrm{Bias}&=\frac{1}{N}\sum_i e_i.
\end{align}
Positive bias denotes under-prediction. Each table states whether it reports a fold mean or a
pooled metric. Fold summaries are the mean $\pm$ sample SD across five folds. Pooled metrics are
calculated once from all out-of-fold predictions. These are not interchangeable.

For the four-survey spatial analyses, pooled $R^2$ 95\% confidence intervals use 1,000 cluster bootstrap resamples of whole compartments. The 2.5th and 97.5th percentiles form the interval. This preserves spatial clustering better than resampling individual rows. The smaller six-survey Q2 results are reported as point estimates; GNNWR is also repeated over three split seeds to test sign stability.

A semivariogram measures how quickly residuals become less similar as distance between plots increases; its range is the distance beyond which they are no longer spatially correlated. Each model's range is fitted separately from its own residuals, searched up to 5000m by default; this window does not widen automatically. One diagnostic only, the six-survey Q2 comparison, was manually rerun with the window widened to 10,000m -- even then several cells remained unresolved at the new cap, reported as provisional, not as a resolved range. Two models' Moran's $I$ are not necessarily tested at the same distance, and any report of a Moran's I value states the distance it was computed at for this reason.

Global Moran's $I$ is calculated from held-out residuals using distance-band weights, with the threshold set by this semivariogram range, and significance assessed via 999 permutations (ranges used are tabulated in the Appendix):

\begin{equation}
I=\frac{N}{W}
\frac{\sum_i\sum_j w_{ij}(e_i-\bar e)(e_j-\bar e)}
{\sum_i(e_i-\bar e)^2}.
\label{eq:morans-i}
\end{equation}
It is used after fitting, not for feature selection.
Three outcomes occur: the range resolves within the search window (Moran's $I$ reported directly); the semivariance stays flat throughout the window (no spatial structure detected, Moran's $I$ not computed); or the semivariance is still rising at the window's edge (Moran's $I$ computed at that capped distance, reported as provisional, not equivalent to a resolved result). The same rule is applied to Q3's bridge check, Q1's residuals and the Elastic Net, XGBoost and GNNWR residuals in Q2.
% moran i 
%   - Weights: binary distance-band (within threshold = neighbour, else not), row-standardised — not inverse-distance, not a kernel
%   - Global Moran's I only, no local variant (LISA) — opposing local clusters could cancel out and go undetected
%   - Detects spatial pattern, not its cause — separate diagnostics (outlier audit, category attribution) needed to explain why
%   - Lower Moran's I is not the same claim as lower residual variance — keep the two metrics distinct in any comparison
%   - Large datasets subsampled to 5,000 rows (fixed seed) before computing, for feasibility

Neural training variation is assessed over five seeds for the selected Q3 models. XGBoost and the
DNN are also compared fold by fold on identical test folds. The physics-weight ablation tests
weights 0, 1 and 2. For PINN-$k$, parameter separability is assessed by the pooled correlation
between learned $y_{\max}$ and $k$; its confidence interval uses a separate 2,000-resample compartment
bootstrap (not the 1,000-resample procedure used for the $R^2$ intervals above). These checks distinguish predictive accuracy from parameter stability.

% Option 2: Short & Crisp (Plain Phrasing)
% \section{Evaluation and robustness}
% \label{sec:eval-robustness}

% Model accuracy is evaluated using RMSE, MAE, $R^2$, and Bias ($e_i = H_i - \hat H_i$):
% \begin{equation}
% \mathrm{RMSE}=\sqrt{\tfrac{1}{N}\textstyle\sum e_i^2}, \quad
% \mathrm{MAE}=\tfrac{1}{N}\textstyle\sum |e_i|, \quad
% R^2=1-\tfrac{\sum e_i^2}{\sum (H_i-\bar H)^2}, \quad
% \mathrm{Bias}=\tfrac{1}{N}\textstyle\sum e_i.
% \label{eq:eval-summary}
% \end{equation}
% Positive bias indicates under-prediction. We report fold-level means ($\pm\text{SD}$) or pooled metrics across all out-of-fold predictions. For the 4-survey cohort, $R^2$ $95\%$ confidence intervals are calculated using 1,000 compartment bootstrap resamples.

% Residual spatial autocorrelation is measured using global Moran's $I$:
% \begin{equation}
% I=\frac{N}{W}\frac{\sum_i\sum_j w_{ij}(e_i-\bar e)(e_j-\bar e)}{\sum_i(e_i-\bar e)^2}.
% \label{eq:morans-i}
% \end{equation}
% Spatial weights ($w_{ij}$) use binary, row-standardised distance bands matched to each model's semivariogram range (up to 5,000\,m; 999 permutations), subsampling 5,000 rows where necessary for computational efficiency. Moran's $I$ is omitted if semivariance is flat, and marked provisional if semivariance rises past 5,000\,m. Distances are reported alongside all values.

% \limit{Global Moran's $I$ detects spatial structure without identifying its cause, and equal-and-opposite local clusters may go undetected. A lower Moran's $I$ reflects less clustering, not lower overall error variance.}

% Neural stability is checked over five seeds. For PINN-$k$, parameter separability is measured by the correlation between predicted $y_{\max}$ and $k$ (2,000 compartment bootstrap resamples).





%=======================================================================================
%=======================================================================================
\chapter{Results} \label{ch:results}
This chapter answers three questions in order. Q1 links long-term growth shifts to environment and management, checking if different models agree. Q2 tests whether local patterns explain what global models miss. Q3 tests if these signals improve top-height prediction and evaluates a physics-informed network.

\section{Q1: What explains a plot's persistent departure from the shared curve?}

Restricting to the 4survey cohort and modelling each plot's mean departure across its
full survey history directly (Table~\ref{tab:results-rq2b}), canopy cover and thinning
history dominate the attribution in every set and every method tested.

\begin{table}[!htbp]
\centering
\scriptsize
\renewcommand{\arraystretch}{1.2}
\begin{tabular}{lrlll}
\toprule
Set & n rows & EN R2 [95\% CI] & XGB R2 [95\% CI] & LMM spatial variance explained \\
\midrule
Set1 (baseline only) & 71{,}766 & 0.220 [0.180, 0.252] & 0.222 [0.183, 0.252] & 0.016$\pm$0.078 \\
Set2 & 71{,}766 & 0.359 [0.310, 0.396] & 0.416 [0.355, 0.461] & 0.204$\pm$0.066 \\
Set3 & 71{,}727 & 0.325 [0.272, 0.365] & 0.375 [0.321, 0.414] & 0.049$\pm$0.048 \\
Set4 & 71{,}330 & 0.358 [0.305, 0.398] & 0.395 [0.338, 0.438] & 0.175$\pm$0.064 \\
\bottomrule
\end{tabular}
\caption{\small{Attribution of the persistent mean curve residual (Equation~\ref{eq:mean-cr-residual}),
4survey only, spatial-block split.}}
\label{tab:results-rq2b} %TODO: CHANGE LABEL, WE ARE USING NEW NAMING CONVENTION
\end{table}

\subsection{Discussion}
Canopy cover carries the largest or near-largest Elastic Net coefficient in every set (Set4: +1.943$\pm$0.034) and ranks first by mean $|$SHAP$|$ in every XGBoost set (1.375--1.442). Because canopy cover is drawn from the same LiDAR pipeline as the height targets themselves, its correlation with all three of this project's own targets was checked directly rather than assumed benign: Pearson correlations of 0.35-0.47, moderate rather than the target restated in different units. Removing canopy cover from Set4 and refitting drops R2 by roughly a third for both models (EN: 0.350 to 0.231, XGBoost: 0.388 to 0.249), a larger effect than its SHAP margin alone (1.2--1.7x over the next variable) would suggest, meaning it carries signal nothing else in the set backfills; the variable most likely to absorb the gap, $\mathit{dist}_{\text{to\_road}}$, does not (its own SHAP value moves slightly the wrong way, 1.003 to 0.869), which rules out simple correlation-driven credit transfer as the mechanism without explaining why canopy cover carries this much unique signal in the first place. The defensible reading is to treat canopy cover and thinning as baseline stand-structure controls, expected to dominate by construction, and to look at the true environmental variables' contribution beyond that baseline as the actual attribution question. 

%TODO: DEFINE ABIOTIC IN THE INTRO OR RELATED WORK.
Among abiotic predictors, only $\mathit{slope\_degrees}$ was stable across both methods, showing consistent positive effects in LMM (+0.894/+0.861) and Elastic Net (+1.230/+0.993) across Sets~3 and 4. In contrast, $\mathit{topex}$ revealed a clear split: LMM found a strong, stable positive link (+1.316/+1.323), whereas Elastic Net yielded weak estimates that flipped signs ($-0.175$ in Set~3 to $+0.063$ in Set~4) with high variance. This reflects a genuine disagreement between models rather than a robust, shared signal. A few reasons below were explored: 

First, multicollinearity was evaluated using variance inflation factors (VIF): all unstable predictors ($\mathit{tas\_mean}$, 3.44; $\mathit{soilgrids\_ph}$, 2.12; $\mathit{topex}$, 2.37) show lower VIF values than the stable $\mathit{chelsa\_bio12\_precip\_mm}$ (3.91), confirming that collinearity does not track instability. 
Second, climatological timing mismatch was examined. While CHELSA features use a static 1981--2010 normal against 2008--2023 survey dates, $\mathit{chelsa\_bio12\_precip\_mm}$ remains stable; conversely, the unstable $\mathit{tas\_mean}$ is derived from HadUK-Grid and averaged over actual survey years. Temporal mismatch therefore fails to separate stable from unstable variables: static features can be stable ($\mathit{chelsa\_bio12\_precip\_mm}$) or unstable ($\mathit{topex}$, $\mathit{soilgrids\_ph}$), while dynamic features can also destabilise. 
Neither multicollinearity nor climate timing explains why $\mathit{soilgrids\_ph}$(a static soil trait that does not change through survey years) is unstable.
%TODO: add this static or changing through survey years in the data table. 

A third possible reason is spatial confounding. Unlike LMM, Elastic Net does not group data by compartment. If a variable mainly changes from one compartment to the next, Elastic Net can easily confuse it with other unmeasured local factors. We checked this by measuring between-compartment variation (ICC). The unstable variables all scored high: $\mathit{topex}$ (0.74), $\mathit{soilgrids\_ph}$ (0.77), and $\mathit{tas\_mean}$ (0.89). However, the stable rainfall variable ($\mathit{chelsa\_bio12\_precip\_mm}$) scored even higher at 0.97. This shows that high spatial variation alone does not make a variable unstable; it makes confounding possible, but is not enough to cause it by itself.

%---
A more unified test explains why Elastic Net destabilises while LMM remains stable. Because both models assume linear relationships, raw non-linearity alone cannot be the cause; the key difference is LMM's compartment random intercept. Mean-centring variables and targets within each compartment isolates this local signal directly:

\begin{itemize}
    \item $\mathit{topex}$ shows a real local signal: its within-compartment correlation rises from 0.09 to 0.21 and becomes monotonic. LMM captures this local effect via random intercepts, whereas Elastic Net pools the data and confuses the signal.
    \item $\mathit{tas\_mean}$ and $\mathit{soilgrids\_ph}$ show almost no local signal: within-compartment correlations collapse from 0.27 to 0.08 and 0.05. Their raw correlations are between-compartment artifacts, leaving no real relationship to estimate reliably.
    \item $\mathit{chelsa\_bio12\_precip\_mm}$ maintains a clean, monotonic local link ($-0.11$ vs.\ $-0.36$ raw), which anchors a stable Elastic Net coefficient despite high spatial variance.
\end{itemize}

Apparent non-linearity in the pooled data is therefore a byproduct of spatial confounding. Variables with genuine local effects stay stable or resolve under random-intercept models, while those driven by between-compartment differences break down (see Appendix~\ref{app:nonlinearity-check}).
%---

All feature sets leave significant spatial autocorrelation in their residuals (distance-band Moran's $I$ at 1385--1739\,m, $p = 0.001$). However, the models remove this clustering at very different rates: Elastic Net drops residual autocorrelation by only 29\% (0.145 $\to$ 0.103), whereas XGBoost reduces it by 63\% (0.146 $\to$ 0.054). This divergence stems from XGBoost's ability to model non-linear interactions automatically \citep{zou2005Regularization_C4_NN}, though its flexibility could also capture spatial noise. Unmodelled clustering persists in both models, pointing to unmeasured environmental factors, resolution limits, or another spatially varying processes \citep{gaspard2019Residual_C5_EVAL}. Set~3 (terrain-only) is the weakest feature set. A category refit ablation on Set~4 shows that soil adds only a small gain ($+0.022$ $R^2$), whereas climate actively reduces accuracy ($-0.056$ $R^2$). Terrain and spatial edge effects drive most of Set~4's advantage. While forest productivity generally relies on mixed biophysical inputs \citep{bontemps2014Predictive_C2_BG}, climate data introduced instability rather than predictive benefit (full ablation in Appendix~\ref{app:category-ablation}). Every environmental set's $R^2$ confidence interval clears the baseline (Set~1) with no overlap, confirming the value of environmental features. However, XGBoost's intervals for Set~2 and Set~4 overlap ([0.355, 0.461] vs.\ [0.338, 0.438]), making their relative ranking inconclusive.

\textbf{For Q1}, stand structure and management account for most of a plot's deviation from the baseline curve. Among abiotic variables, only slope is consistently stable. Instability in other features depends on whether a genuine local signal survives between-compartment pooling: LMM recovers a real local relationship for $\mathit{topex}$ that Elastic Net obscures, whereas $\mathit{tas\_mean}$ and $\mathit{soilgrids\_ph}$ are largely between-compartment artifacts. Stable rainfall ($\mathit{chelsa\_bio12\_precip\_mm}$) retains a robust local signal despite high spatial variance. At the category level, soil adds a slight gain while climate harms performance. Because significant spatial autocorrelation persists across both models (Moran's $I$ falls 29\% in EN vs.\ 63\% in XGBoost), global environmental features fail to capture all spatial structure, directly justifying the spatially varying models evaluated in Q2.


\fig{\textbf{Q1 -- coefficient/Moran's I/residual-map composite.} Three panels.
Panel A (attribution): coefficient forest plot, one row per variable (\texttt{CanopyCover},
both thinning columns, \texttt{slope\_degrees}, \texttt{topex}, the climate/soil
variables), point = mean, whisker = fold SD, LMM and EN as two colours/markers side by
side, Set3 and Set4 as two sub-columns -- the visual payoff of the slope-stable/
topex-unstable contrast, otherwise the reader has to hold four numbers in their head.
Panel B (spatial clustering responds unevenly to environmental data): Moran's I by set
(Set1--Set4), EN and XGBoost as two lines -- EN nearly flat, XGBoost dropping
substantially, the divergence between the two lines is the finding, not a shared trend.
Panel C (where the leftover residual sits): a map of Set4/4survey's model residual,
coloured by magnitude/sign -- this is what motivates Q2, earning its panel on
structural grounds, not decoration. Establishes: none of this (coefficients, the
Moran's I trend, the spatial residual pattern) exists in any table in this chapter, all
of it is currently text-only. Answers Q1 items 1--3.}


%maybe when we comapre differencce ebtween 4 survey and 6 survey, we can say that as 6 survey is morehomogenous, we know it has a sifferent CR fit, then (((this oudlbe a reaosn or justifcation) for something)) 

\textbf{For Q1}, stand structure and management explain most of the attributable growth deviations. Among environmental factors, only slope shows a consistent relationship across model families. A large share of spatial variation remains unexplained by any tested environmental set, directly motivating the use of a spatially varying model in Q2.

%=======================================================================================
\section{Q2: Spatial attribution}
Q1 showed that global models leave notable spatial clustering in the residuals, with XGBoost removing far more autocorrelation than Elastic Net with added environmental data. Q2 directly tests whether letting these relationships vary spatially captures what global models miss (Table~\ref{tab:results-rq3}).
\begin{table}[!htbp]
\centering
\scriptsize
\renewcommand{\arraystretch}{1.2}
\begin{tabular}{lllllll}
\toprule
Set & EN 4survey [CI] & EN 6survey & XGB 4survey [CI] & XGB 6survey & GNNWR 4survey [CI] & GNNWR 6survey \\
\midrule
Set2 & 0.286 [0.241, 0.325] & 0.057 & 0.266 [0.214, 0.308] & 0.031 & 0.320 [0.270, 0.369] & $-0.044$ / 0.014$\pm$0.053 \\
Set3 & 0.274 [0.223, 0.316] & 0.031 & 0.281 [0.235, 0.320] & 0.086 & 0.319 [0.263, 0.365] & 0.058 / 0.054$\pm$0.055 \\
Set4 & 0.240 [0.192, 0.286] & 0.056 & 0.250 [0.201, 0.295] & 0.015 & 0.294 [0.236, 0.347] & $-0.075$ / 0.002$\pm$0.068 \\
\bottomrule
\end{tabular}
\caption{\small{Attribution of the plot-specific asymptote deviation
(Equation~\ref{eq:local-ymax-diff}), spatial-block split. 6survey GNNWR column shows
seed-42 point estimate / three-seed mean$\pm$SD.}}
\label{tab:results-rq3} %TODO: rename table label refirng to q2 not rq3
\end{table}

GNNWR achieves higher point estimates than both global baselines across all 4-survey sets, though bootstrap confidence intervals overlap in every set (Set~4: GNNWR [0.236, 0.347] vs.\ EN [0.192, 0.286]). Residual autocorrelation supports this directional advantage: GNNWR reduces Moran's $I$ beyond the best global baseline across all 4-survey sets (Set~2: $-32.0\%$; Set~3: $-24.7\%$; Set~4: $-23.7\%$; $p = 0.001$), confirming it models real spatial structure. On the 6-survey set, the diagnostic is uninformative because semivariograms fail to resolve, leaving no clear spatial signal to evaluate. While spatial models are sensitive to sample density \cite{brunsdon1996Geographically_C6_GWR} and the 6-survey cohort spans far fewer compartments (47 vs.\ 231), downsampling 4-survey reference points down to 7,500 failed to trigger a similar collapse ($R^2$ stayed positive at 0.12--0.34). Having fewer points alone does not explain the collapse, the ablation reduced point counts evenly across all compartments, whereas the 6-survey dataset is restricted to far fewer compartments to begin with.

% Option 1: Direct & Academic (Recommended)
The same $\sim$10 plots consistently produce the largest residuals across nearly every feature set and cohort, ruling out poor feature selection as the cause. This recurrence holds across independent model families: flagged plots fall into the worst 1\% and worst 5\% absolute residual buckets at roughly 36 times the rate expected by chance across Elastic Net (35.7\%/59.0\%), XGBoost (38.0\%/61.3\%), and GNNWR (37.6\%/57.5\%).  None of these 10 plots trigger standard data-cleaning filters for disturbance, measurement anomalies, or trajectory instability. Instead, inspecting their trajectories reveals two distinct causes. Half display smooth, consistent growth whose large deviations reflect a mismatch with the baseline yield-class lookup rather than faulty raw heights. The remaining half show sub-threshold curve instability across survey years, producing implausible fitted asymptotes (up to 116\,m despite raw heights of 38.9--46.8\,m). Tracing these asymptotes back to the fitting procedure confirms the mechanism: because the closed-form curve holds the growth-rate parameter fixed to the assigned yield class, a plot growing slower or flatter than assumed is forced into an inflated asymptote. This issue stems directly from target construction rather than raw data quality, meaning simple height filters cannot detect it. Although GNNWR shows 2.7 times higher error on these flagged plots (21.12\,m vs.\ 7.70\,m), its local coefficients remain identical to control plots. This confirms that the large errors are target-construction artifacts rather than real environmental or spatial anomalies.
%TODO add these rules to the methodlogy

These flagged cases extend beyond a few isolated outliers: a Tukey-fence check on fitted asymptotes flags 0.47\% of the 4-survey cohort (266 plots) and 5.32\% of the 6-survey cohort (709 plots). These plots cluster tightly within their compartments rather than dispersing evenly, pointing to localised conditions rather than compartment-wide labelling errors. However, removing flagged plots does not improve performance: pooled $R^2$ barely shifts in the 4-survey set ($-0.002$) and worsens in the 6-survey set (Elastic Net $-0.027$, XGBoost $-0.037$). Dropping them removes usable signal rather than cleaning noise. None of the 266 plots trigger flags for clearfelling or measurement error, ruling out abrupt catastrophes while leaving chronic, sub-threshold effects possible.

Testing whether the fixed-$k$ mechanism applies across the broader flagged set reveals two distinct groups rather than a single artifact:
\begin{itemize}
    \item \textbf{Artifact-driven subset (71.6\%, 179/250):} Plots growing slower than their assigned yield-class benchmark show the direct signature of the fixed-$k$ constraint, where slow growth forces an inflated asymptote. These plots are concentrated on exposed terrain (median $\mathit{topex}$ $-5.48$) with modest evidence of chronic wind exposure ($p = 0.0007$).
    \item \textbf{Fast-growing subset (28.4\%, 71/250):} Plots whose observed growth matches or exceeds the benchmark curve do not fit the fixed-$k$ artifact explanation. These plots sit in sheltered terrain (median $\mathit{topex}$ $+2.58$, $p = 0.0002$) and have more recent, frequent thinning ($p = 0.051$), suggesting genuinely accelerated growth under favorable local conditions.
\end{itemize}
The fixed-$k$ curve-fitting artifact explains most flagged asymptotes in exposed terrain, whereas extreme asymptotes in sheltered areas likely reflect true rapid growth rather than fitting errors (full diagnostics are detailed in Appendix~\ref{app:flagged-plot-split}).

Variable importance rankings vary by cohort: canopy cover ranks first in the 4-survey cohort across three converging attribution metrics (Elastic Net coefficients, XGBoost gain importance, and SHAP values), whereas XGBoost disagrees with Elastic Net and GNNWR on the 6-survey cohort. However, because XGBoost achieves near-zero or negative $R^2$ on the 6-survey data (Table~\ref{tab:results-rq3}), its attribution rankings there warrant less confidence. Given GNNWR's established instability on the 6-survey subset, the 4-survey cohort provides the more reliable baseline for addressing Q2's spatial questions. Grouping Set~4 predictors into categories confirms that stand structure dominates overall importance (roughly $3\times$ the next-largest category). Among abiotic groups, terrain provides the largest contribution, while soil is indistinguishable from noise, aligning with the individual variable results above (category-level Moran's $I$ diagnostics, including an adjustment for a bandwidth artifact in the terrain check, are reported in Appendix~\ref{chap:appendix}).

% Option 2: Short & Crisp (Simpler Phrasing)
\textbf{For Q2}, results support the view that environmental relationships vary spatially, provided there are enough nearby training plots for the spatial-weighting network to learn reliable local patterns. Because having fewer training plots alone does not cause the 6-survey collapse, the difference between cohorts is not just an artifact of total sample size. Although GNNWR's $R^2$ edge overlaps with global confidence intervals, its greater reduction in residual Moran's $I$ confirms it captures real spatial patterns. Furthermore, three distinct attribution methods agree that most extreme deviations arise from target construction under a fixed growth rate. GNNWR finds normal local coefficients at these locations, showing that the largest errors are primarily mathematical artifacts rather than localized environmental signals. Crucially, this mechanism accounts for most, but not all, outliers: unexplained cases sit in sheltered terrain, while the artifact-linked majority correlates with chronic wind exposure. Extreme deviation values must therefore be interpreted cautiously, as they reflect a mix of curve-fitting artifacts and genuine biological growth.


\fig{\textbf{Q2, Figure 1 -- GNNWR edge/spatial-structure/local-coefficient
composite.} Three panels. Panel A: pooled R2 by set and model (EN/XGBoost/GNNWR),
point+CI, 4survey only. Panel B: residual Moran's I by set/model, 4survey only (mark
6survey's cells as unresolved/not-significant rather than plotting a misleading point
value) -- the independent corroboration for Panel A's edge, and the reason 6survey is
absent rather than shown as a contradictory result. Panel C: a spatial map of GNNWR's
own per-plot local
coefficient for \texttt{CanopyCover} (its top-ranked variable), coloured by
coefficient value -- the only way to show what "spatially-varying coefficient" means
concretely, no prose substitutes for seeing it vary across Aberfoyle. Establishes:
Panel B is a diagnostic that exists nowhere else visually; Panel C is the one thing in
the whole chapter that only GNNWR can produce. Answers Q2 items 2--3.}

\fig{\textbf{Q2, Figure 2 -- deviation map/archetype map/trajectory composite,
including the 266-plot pattern.} Three panels. Panel A: a map of
\texttt{local\_y\_max\_difference} across 4survey, coloured by deviation
magnitude/sign, with the 266 Tukey-fence-flagged plots marked with a distinct
outline/marker -- the direct visual for "the pattern scales to hundreds of plots,
clustered by compartment" (item 6), and shows at a glance whether it clusters
spatially. Panel B: the compartment archetype map (the six classes from item 9), same
geography, so a reader can cross-reference Panel A's raw deviation against Panel B's
classification directly. Panel C: small-multiples trajectories, one per archetype,
observed height vs.\ age (grey), the fitted local curve, and the yldc benchmark curve
overlaid -- the only way to make "smooth-but-offset" vs.\ "moderately unstable"
concrete rather than a label. Establishes: this figure is what lets the outlier-plot
prose (items 4--9, currently the densest part of the chapter) shorten, rather than the
other way round. Answers Q2 items 4--6, 9.}




%=======================================================================================
\section{Q3: Prediction and physics guidance}
\label{sec:results-q3}
Q1 and Q2 established that a plot's persistent deviation from the shared growth curve carries a real, if partial and spatially variable, environmental attribution signal. This supplementary section returns to the dissertation's primary question: does this signal translate into a predictive advantage for top height itself, and can deep neural networks exploit it more effectively than tree-based ensembles like XGBoost? Across all models, cohorts, and spatial-block splits, tuned XGBoost delivers the most accurate top-height predictions. CR loss guidance does not help the neural network match XGBoost's accuracy at any tested weight ($\lambda$), but it still successfully separates the asymptote ($y_{\max}$) from the growth rate ($k$).

\subsection{Results: top-height prediction:}
\begin{table}[!htbp]
\centering
\scriptsize
\setlength{\tabcolsep}{4pt}
\renewcommand{\arraystretch}{1.2}
\begin{tabular}{lrrrrrr}
\toprule
& \multicolumn{3}{c}{Four-survey} & \multicolumn{3}{c}{Six-survey} \\
\cmidrule(lr){2-4} \cmidrule(lr){5-7}
Model & RMSE${\downarrow}$ & MAE${\downarrow}$ & $R^2{\uparrow}$ & RMSE${\downarrow}$ & MAE${\downarrow}$ & $R^2{\uparrow}$ \\
\midrule
Linear & $5.168{\pm}0.361$ & $4.056{\pm}0.343$ & $0.580{\pm}0.026$ & $4.877{\pm}0.643$ & $3.619{\pm}0.426$ & $0.509{\pm}0.067$ \\
RF & $4.794{\pm}0.187$ & $3.611{\pm}0.142$ & $0.638{\pm}0.011$ & $4.062{\pm}0.249$ & $3.010{\pm}0.185$ & $0.656{\pm}0.046$ \\
XGBoost (tuned) & $\mathbf{4.548{\pm}0.137}$ & $\mathbf{3.431{\pm}0.123}$ & $\mathbf{0.674{\pm}0.009}$ & $\mathbf{3.656{\pm}0.232}$ & $\mathbf{2.669{\pm}0.183}$ & $\mathbf{0.722{\pm}0.031}$ \\
DNN & $4.681{\pm}0.230$ & $3.533{\pm}0.201$ & $0.655{\pm}0.016$ & $3.903{\pm}0.316$ & $2.867{\pm}0.275$ & $0.684{\pm}0.031$ \\
PINN$^*$ & $5.206{\pm}0.399$ & $4.063{\pm}0.407$ & $0.573{\pm}0.036$ & $3.886{\pm}0.290$ & $2.819{\pm}0.196$ & $0.686{\pm}0.038$ \\
PINN-$k^*$ & $5.197{\pm}0.402$ & $4.046{\pm}0.421$ & $0.575{\pm}0.037$ & $3.895{\pm}0.282$ & $2.816{\pm}0.196$ & $0.684{\pm}0.038$ \\
\bottomrule
\end{tabular}
\caption{\small{Model comparison for top-height prediction, Set3, spatial\_block\_kfold,
per-fold mean $\pm$ SD across five folds. 4survey n=232{,}292 row-years (58{,}073
plots); 6survey n=82{,}614 row-years (13{,}769 plots). Table shows pooled mean $\pm$ SD across all folds, and margin reflects the full split rather than single folds.
$^*$PINN/PINN-$k$: the implemented forward pass does not carry environmental
information to the prediction itself, only to auxiliary loss terms -- see
``PINN: what it did and did not test'' below. Results are for the architecture
as implemented, not as intended.}}
\label{tab:results-rq1}
\end{table}

\subsection{Discussion}
This section evaluates the direct comparison between XGBoost and the standard DNN, where environmental features feed directly into top-height predictions. Across both cohorts, XGBoost consistently outperforms the DNN. Evaluating individual spatial cross-validation folds confirms this is not an artifact of pooling: XGBoost achieves higher accuracy in 9 of the 10 folds (all 5 folds on the 6-survey cohort, and 4 of 5 on the 4-survey cohort, trailing the remaining fold by only $-0.007$ $R^2$). Four candidate explanations account for this performance gap:

\begin{itemize}
\item \textbf{Feature interactions (insufficient explanation):} XGBoost allocates 27\% of its SHAP prediction weight to interactions (primarily Age crossed with terrain and stand structure). However, DNN already captures 80\% of the nonlinear gain above the linear baseline ($R^2$ 0.580 $\to$ 0.655 out of a possible 0.674; Table~\ref{tab:results-rq1}), and Random Forest underperforms the DNN ($R^2 = 0.638$). Interaction capacity alone therefore does not explain the margin.
\item \textbf{Compartment sample size (ruled out):} Subsampling the 4-survey cohort down to 48 compartments (matching the 6-survey count) did not widen XGBoost's lead; the $R^2$ difference fluctuated between $-0.085$ and $+0.029$ across random subsets, compared to $+0.024$ on full data.
\item \textbf{Architecture-size under-tuning (ruled out):} Sweeps across five network sizes showed no consistent gains for DNN.
\item \textbf{Terrain extrapolation in held-out compartments (ruled out):} On the 4.1\% of 4-survey test plots with out-of-range terrain values, XGBoost degraded more relative to its in-range error than DNN did (error ratio 1.312 vs.\ 1.243; correlation with out-of-range distance $\rho=0.074$ vs.\ $0.059$).
\item \textbf{Training hyperparameter under-tuning (ruled out):} A 48-configuration sweep over dropout and learning rate yielded no meaningful gains for neural models under 5-fold cross-validation (within $\pm0.008$ $R^2$ of default baselines). Similarly, XGBoost's 27-configuration tuning shifted $R^2$ by only 0.002-0.004, confirming that neither model family had substantial untapped hyperparameter headroom.
\end{itemize}

The remaining theory is that, boosted-tree ensemble like XGBoost are better suited to this tabular task than the neural network architecture, which aligns with standard literature on gradient-boosted trees outperforming deep learning on heterogeneous tabular data \cite{shwartzziv2021Tabular_C4_NN}.

\paragraph*{Does environmental conditioning correct the curve at row level?}
Evaluating row-level corrections against the fold-matched Chapman-Richards baseline (Equation~\ref{eq:residual-reduction}) shows that environmental conditioning reduces error unevenly. The gains concentrate heavily in the worst-fitting quartile (reducing error by up to 5.08\,m in the 4-survey cohort), whereas the best-fitting quartile degrades slightly. A permutation control with shuffled environmental features reveals that the worst-quartile improvement shrinks by 18--27\% rather than disappearing entirely, while the best-quartile penalty persists at a similar magnitude. This indicates that correcting large errors relies partly on genuine environmental signals, whereas the slight cost in the best-fitting quartile reflects a general flexibility penalty of more complex models rather than an issue specific to environmental conditioning.

\paragraph*{Answer to Q3.} Tuned XGBoost gives the strongest predictions in both cohorts; the DNN stays competitive; CR guidance does not improve accuracy for the architecture actually tested.

% Option 1: Direct & Academic (Recommended)
\paragraph*{PINN: theoretical value versus implementation scope.}
\label{sec:rq1-pinn-limitation}
In theory, physics-informed neural networks (PINNs) offer a compelling balance between empirical machine learning and domain knowledge. By embedding the CR growth equation into the loss formulation, a PINN provides explicit physical interpretability: it can estimate plot-specific biological parameters, namely the asymptotic maximum height ($y_{\max}$) and the intrinsic growth rate ($k$), allowing the resulting site-specific curves and equations to be inspected and understood. In this study, however, PINN and $\text{PINN}_k$ trailed tuned XGBoost and the standard DNN across all evaluated physics loss weights ($\lambda$), with higher penalties consistently degrading top-height predictive accuracy (e.g., removing the physics penalty entirely yielded a $+0.05$ $R^2$ improvement on the 4-survey cohort).

This performance gap reflects an implementation constraint rather than evidence against the theory of physics-informed modelling. In the evaluated architecture, the per-plot adjusted asymptote and growth rate were constrained only via the auxiliary physics and trajectory loss terms, without being carried directly through the forward pass of the main prediction trunk. Two diagnostics confirm this reading:
A random (non-spatial) split inflated the standard DNN's $R^2$ from $0.634$ to $0.831$ due to spatial data leakage, whereas the PINN remained essentially unchanged ($0.584$ vs.\ $0.590$), confirming that environmental features never directly reached the PINN trunk to leak. Setting the physics loss weight to zero eliminated most of the DNN-PINN performance gap on the 4-survey dataset (reducing the margin from $0.080$ to $0.026$).

While this structural disconnect prevented the model from fully realising its theoretical interpretability and predictive synergy within this dissertation's scope, an architectural fix (explicitly injecting the evaluated Chapman-Richards curve into the output layer) has been designed and documented in Appendix~\ref{chap:appendix}.

\limit{These results reflect the evaluated architecture rather than the fully realized theoretical PINN; a corrected forward pass could yield different predictive and interpretive behaviour.}









%=======================================================================================
%=======================================================================================
\chapter{Conclusion and Future Work} \label{ch:conc}

%=======================================================================================
%=======================================================================================
\bibliographystyle{plain}
\bibliography{mybibfile}
%=======================================================================================

\appendix
\chapter{Appendix}  \label{chap:appendix} 
% \TODO: create appendix secions and labels for each section. 



%=======================================================================================
\end{document}
%=======================================================================================

